% Front-matter highlights. One line per load-bearing result, each pointing at
% the section that proves it. Uses only packages already in preamble.tex.
\section*{Highlights}
\label{sec:highlights}
\addcontentsline{toc}{section}{Highlights}

\begin{itemize}[leftmargin=1.6em,itemsep=0.35em]
  \item \textbf{The rotation is structural, not incidental.} The controller of
        \cite{kou2022cooperative}, Eq.~(7), constrains only the \emph{average} position error and the
        \emph{average} velocity, leaving $2(N-1)$ circulation degrees of freedom unconstrained by the
        average dynamics (\S\ref{sec:avg}).

  \item \textbf{But those degrees of freedom are not undamped.} The isotropic gain $k_1$ damps every
        direction. What singles out rotation is a \emph{cancellation}: at radial balance the tangential
        stiffness of the repulsion exactly offsets $k_1$, leaving one neutral direction---the rotational
        phase---with every other mode strictly decaying (\S\ref{sec:corotating}).

  \item \textbf{Rotation is the only alternative to collapse.} The observer
        $\dot{\tvell}_i=-k_2\txi_i$ has no damping term, so a persistent nonzero position error is
        integrated without bound. A static formation would need a linearly growing velocity estimate
        balanced by static forces, which is impossible; the only escape is to rotate
        (Remark~\ref{rem:trichotomy}).

  \item \textbf{The angular speed is forced, exactly.} Any nondegenerate rigid rotation satisfies
        $\abs{\omega}=\sqrt{k_2}$, independently of the radii and of the shape, because the net internal
        torque of central pairwise forces vanishes. The hypothesis $\omega\ne0$ is \emph{not} needed---it
        follows from nondegeneracy (Proposition~\ref{prop:omega}).

  \item \textbf{The rotating formation need not be a regular polygon.} Initialized on a regular hexagon,
        the steady state settles into an interleaved two-triangle pattern with alternating radii, yet all
        six vehicles share the same $\sqrt{k_2}$: the zero-torque identity eliminates every radius before
        $\omega$ is determined (\S\ref{sec:simulation}).

  \item \textbf{The centroid sits exactly on the target.} The rigid-rotation ansatz forces
        $\bx^0=0$, so this is a conclusion rather than a modeling assumption
        (Remark~\ref{rem:centroid}).

  \item \textbf{Local orbital stability holds; a global proof is obstructed.} The zero-mean shape
        subsystem $\ddot{\txi}^s+B(\txi^s)\dot{\txi}^s+k_2\txi^s=0$, with $B=k_1I-\nabla^2P$, has energy
        $E$ satisfying $\dot E=-\dot{\txi}^{s\mathsf T}B\dot{\txi}^s$. Rotation covariance makes the
        rotational direction a null mode of $B$ at radial balance, and the co-rotating spectrum confirms
        stability for the paper's equilibrium---but $B$ is sign-indefinite away from equilibrium, so the
        energy argument does not globalize (\S\ref{sec:corotating}).

  \item \textbf{The individual position errors are bounded.} Only the average is forced to decay, yet
        the raw errors $\txi_i=x_i-x_0$ cannot diverge: the functional
        $W=\frac12\sum_i\norm{\tvell_i}^2+\frac{k_2}{2}\sum_i\norm{\txi_i}^2$ satisfies
        $\dot W=\frac{k_2}{2}\sum_{i\ne j}\alpha(\norm{x_{ij}})\norm{x_{ij}}-k_1k_2\sum_i\norm{\txi_i}^2$,
        giving an explicit bound whenever gaps stay off the barrier, while near the barrier the
        proximity components are isolated clusters whose centroids obey the same damped oscillator as
        the average (\S\ref{sec:bounded}).

  \item \textbf{Speed is determined; sign is not.} The sign of $\omega$ equals the sign of the directed
        circulation $L=\sum_i\txi_i\times\tvell_i$, and any nonzero initial circulation selects the
        rotation sense. The exactly symmetric case is degenerate and, in the runs performed, does not
        reach a rigid rotation within the horizon tested (\S\ref{sec:direction}).

  \item \textbf{The second controller stops the rotation for a specific algebraic reason.} Adding
        $\phi_i$ to the observer makes $\partial V/\partial x_i$ the exact negative of the observer's
        right-hand side, turning $\dot V$ into a sum of negative squares and dissipating \emph{every}
        individual velocity error, not just the average (\S\ref{sec:second}).

  \item \textbf{The mechanism survives to acceleration inputs.} With a uniformly accelerating target the
        same torque balance gives $\abs{\omega}=\sqrt{k_3/k_2}$, with regular-hexagon radius solving
        $\alpha(r)=(k_1-k_3/k_2)r$, which exists when $k_1k_2>k_3$ (\S\ref{sec:double}).

  \item \textbf{Three branches, one open question.} Generic initializations rotate; symmetric ones
        jam onto the collision barrier, with gaps closing to $d$ at rate $1/t$ while observer states
        diverge linearly; and at sufficiently large initial radii the formation enters a breathing
        limit cycle with periodic shape oscillations. In the collinear case $\nabla^2P\preceq0$ makes
        $B\succeq k_1I\succ0$ and the jammed branch is provable; whether the trichotomy is exhaustive
        and whether it can be refined into a theorem remains
        Conjecture~\ref{conj:dist}.
\end{itemize}

\noindent
All numerical experiments, integration settings, and reproduction scripts are collected in the
standalone companion \emph{Numerical Experiments for ``Why Does the First Fencing Controller Induce
Rotation About the Target?''}; references of the form \S\ref{S-sec:sim_omega} point into it.
